% --------------------------------------------------------------------
%  PACKAGES
% --------------------------------------------------------------------
%  The union of what the books need. A course that needs something
%  extra loads it in its own file, under the heading "packages only
%  this book needs". Pseudocode is the one optional package the shared
%  file knows about: \utpusealgorithms in listings.tex loads and styles it.
%
%  The tikz and pgfplots library lists are the union across all books,
%  so a diagram or chart copied from one course compiles in another.

% Register the Source Sans map where MiKTeX has not registered it globally.
\ifdefined\pdfmapfile
  \pdfmapfile{+SourceSansThree.map}
\fi
\usepackage[T1]{fontenc}
\IfFileExists{glyphtounicode.tex}{\input{glyphtounicode}\pdfgentounicode=1}{}
\usepackage[american]{babel}
\IfFileExists{sourcesanspro.sty}
  {\usepackage[default]{sourcesanspro}}
  {\usepackage[scaled]{helvet}\renewcommand{\familydefault}{\sfdefault}}
\DeclareFontShape{T1}{cmss}{b}{n}{<->ssub*cmss/bx/n}{}
\DeclareFontShape{TS1}{cmss}{b}{n}{<->ssub*cmss/bx/n}{}
% Monospace. Without this the tt default stays Computer Modern Typewriter,
% a serif face that clashes with the sans body text in every listing and
% every \cmd. Source Code Pro is the companion to Source Sans Pro; it is
% scaled down because mono digits otherwise sit large next to the sans.
%
% Inconsolata, chosen with \def\utpmonofont{inconsolata}, distinguishes
% 0/O and 1/l/I at a glance, which matters when a reader copies an
% interface name or an address out of a transcript. Only one mono package
% is ever loaded: a second would win the family but inherit the first's
% series names, and every \ttfamily would then resolve by substitution.
\def\utpmonoinconsolata{inconsolata}
\ifx\utpmonofont\utpmonoinconsolata
  \IfFileExists{inconsolata.sty}{\usepackage[varqu,varl,scaled=1.02]{inconsolata}}{}
\else
  \IfFileExists{sourcecodepro.sty}
    {\usepackage[scaled=0.88]{sourcecodepro}}
    {\usepackage[scaled=0.86]{beramono}}
  % Source Sans names its normal weight "regular", but the Computer Modern
  % fallbacks call it "m". Map one onto the other so a \ttfamily reached
  % while the regular series is active resolves instead of warning.
  \DeclareFontShape{T1}{cmtt}{regular}{n}{<->ssub*cmtt/m/n}{}
  \DeclareFontShape{TS1}{cmtt}{regular}{n}{<->ssub*cmtt/m/n}{}
  \DeclareFontShape{T1}{cmtt}{regular}{it}{<->ssub*cmtt/m/it}{}
\fi
\usepackage[margin=21mm,headheight=22pt]{geometry}
\usepackage{amsmath,amssymb,mathtools,bm}
\usepackage{microtype}
\usepackage{enumitem}
\usepackage{booktabs}
\usepackage{tabularx}
\usepackage{array}
\usepackage{ragged2e}
\usepackage{etoolbox}
\usepackage{titlesec}
\usepackage{tocloft}
\usepackage{fancyhdr}
\usepackage[table]{xcolor}
\usepackage{tcolorbox}
\tcbuselibrary{breakable,skins}
\usepackage[hidelinks]{hyperref}
\usepackage{bookmark}
% Document properties, so the PDF identifies itself in a reader's
% title bar, in a library catalogue, and in search results. Set from
% the same configuration macros the cover uses.
\hypersetup{
  pdftitle={\utpcoursetitle},
  pdfauthor={\utpauthor},
  pdfsubject={Universiti Teknologi PETRONAS course notes and laboratory guide},
  pdfkeywords={UTP, course notes, laboratory manual, \utpcoursetitle},
  pdfdisplaydoctitle=true,
  bookmarksnumbered=true,
  bookmarksopenlevel=1
}
\usepackage{listings}
\usepackage{graphicx}
\usepackage{caption}
\usepackage{tikz}
\usetikzlibrary{angles,arrows.meta,backgrounds,calc,chains,decorations.pathmorphing,decorations.pathreplacing,fit,matrix,patterns,positioning,quotes,shadows.blur,shapes.geometric,shapes.misc,shapes.symbols}
\usepackage{pgfplots}
% statistics gives boxplot prepared; groupplots gives multi-panel charts.
\usepgfplotslibrary{statistics}
\usepgfplotslibrary{groupplots}

% Images: any folders the book names first, then its own figures/
% folder, then the shared logos used by the cover.
\expandafter\graphicspath\expandafter{\utpextragraphicspath{figures/}{../.shared/logos/}}
